\chapter{KẾT LUẬN VÀ HƯỚNG PHÁT TRIỂN}
\label{Chapter5}

\section{Kết luận chung}
Nghiên cứu đã xây dựng và kiểm chứng thành công một khung phương pháp toàn diện, bao quát trọn vẹn các bước từ thu thập dữ liệu, tiền xử lý, gán nhãn cho đến tinh chỉnh mô hình, nhằm giải quyết bài toán phân loại cảm xúc vi mô trên ngữ liệu mạng xã hội tiếng Việt. Giá trị cốt lõi của khung phương pháp này là khả năng xử lý triệt để những thách thức đặc thù của dữ liệu thực tế như văn bản thiếu ngữ cảnh, độ nhiễu cao và sự mất cân bằng trong phân phối. Việc chuẩn hóa quy trình không chỉ giúp tối ưu nguồn lực tính toán, mà còn tạo dựng một hệ thống linh hoạt, có khả năng mở rộng cho các bài toán phân tích ngữ nghĩa phức tạp hơn. Trên cơ sở lý thuyết cùng các minh chứng thực nghiệm định lượng, nghiên cứu rút ra những kết luận chính như sau.

\textbf{Thứ nhất, về tính thực tiễn của quy trình gán nhãn bán tự động:} Việc kết hợp tập luật từ vựng (rule-based) với các Mô hình Ngôn ngữ Lớn (LLMs) dựa trên tập hạt giống (seed data) đã cho thấy hiệu quả vượt trội, giải quyết thành công bài toán thời gian và chi phí vốn là điểm yếu của phương pháp gán nhãn chuyên gia truyền thống khi xử lý dữ liệu quy mô lớn. Đáng chú ý, phân phối nhãn thu được phản ánh trung thực cấu trúc đuôi dài tự nhiên của ngôn ngữ mạng xã hội, vừa đảm bảo độ chính xác vừa hạn chế tối đa tình trạng gán nhãn sai cho các nhóm cảm xúc thiểu số.

\textbf{Thứ hai, về năng lực mô hình và kỹ thuật xử lý mất cân bằng dữ liệu:} Kết quả thực nghiệm cho thấy ưu thế rõ rệt của kiến trúc \textbf{mModernBERT}, nhờ những cải tiến trong cơ chế Attention và khả năng mở rộng ngữ cảnh, mô hình đạt chỉ số F1-Macro vượt quá $88\%$. Thành tựu này phần lớn đến từ chiến lược xử lý mất cân bằng dữ liệu: kết hợp kỹ thuật giảm mẫu đa số với hàm mất mát hỗn hợp \textbf{Mixture Loss} (gồm CB-Loss và GHM-Loss), giúp điều chỉnh hiệu quả trọng số học tập của mô hình. Nhờ đó, hệ thống không bị chi phối bởi các nhãn phổ biến mà vẫn giữ được độ nhạy cao với những cảm xúc tiêu cực vi mô hiếm gặp như sợ hãi (\texttt{fear}), hối tiếc (\texttt{remorse}) hay thương tiếc (\texttt{grief}).

\textbf{Tóm lại,} sự kết hợp chặt chẽ giữa quy trình tự động hóa dữ liệu và các kỹ thuật điều chỉnh mô hình chuyên sâu đã chứng minh được tính chính xác, độ tin cậy và toàn diện của khung phương pháp đề xuất. Nghiên cứu không chỉ mang lại một giải pháp học máy tối ưu cho việc Phân tích Cảm xúc tiếng Việt, mà còn khẳng định ứng dụng vào thực tiễn mạnh mẽ, tạo nền tảng để triển khai các hệ thống lắng nghe xã hội (social listening) và giám sát dư luận tự động với độ chính xác cao.

\section{Hướng phát triển}

Mặc dù khung phương pháp đề xuất đã đạt nhiều kết quả khả quan trong thực nghiệm và cho thấy được tính khả thi, một số giới hạn vẫn cần được nhìn nhận. Trên cơ sở đó, các hướng nghiên cứu tiếp theo sẽ tập trung hoàn thiện và mở rộng khả năng ứng dụng của hệ thống.

Hiện nay, nghiên cứu vẫn chủ yếu khai thác đặc trưng ngữ nghĩa từ dữ liệu văn bản (text-based), trong khi trên thực tế, cảm xúc người dùng ở các nền tảng mạng xã hội thường chịu ảnh hưởng từ bối cảnh video, hình ảnh hay âm thanh đi kèm. Việc thiếu hụt dữ liệu đa phương thức vì vậy có thể dẫn đến sai lệch khi phân tích những bình luận mang tính châm biếm hoặc sử dụng tiếng lóng phụ thuộc ngữ cảnh.

Bên cạnh đó, dù quy trình tạo nhãn giả (pseudo-labeling) đã được kiểm soát chặt chẽ, việc phụ thuộc vào đầu ra của Mô hình Ngôn ngữ Lớn vẫn tiềm ẩn rủi ro sai sót đối với những trạng thái cảm xúc nằm ở vùng ranh giới.

Ngoài ra, các kiến trúc Transformer cỡ lớn như mModernBERT đòi hỏi tài nguyên phần cứng đáng kể, gây khó khăn về độ trễ khi cần triển khai hệ thống phân tích thời gian thực ở quy mô lớn.

Nhằm khắc phục những hạn chế vừa nêu, các hướng phát triển tiếp theo sẽ tập trung vào:

- Bổ sung thêm luồng dữ liệu âm thanh và video bên cạnh văn bản, nhằm xây dựng một không gian biểu diễn ngữ nghĩa toàn diện hơn, qua đó hỗ trợ tốt hơn cho bài toán nhận diện cảm xúc châm biếm cũng như các bối cảnh còn ẩn khuất.

- Tích hợp cơ chế Học chủ động (Active Learning) vào khung phương pháp, cho phép hệ thống tự động chọn lọc những mẫu dữ liệu mà mô hình dự đoán với độ tin cậy thấp nhất để chuyển cho chuyên gia thẩm định, trước khi cập nhật trở lại tập huấn luyện nhằm cải thiện chất lượng dữ liệu.

- Áp dụng các kỹ thuật như lượng tử hóa (Quantization) hoặc chưng cất tri thức (Knowledge Distillation) nhằm thu gọn kiến trúc mModernBERT thành phiên bản nhẹ hơn, vừa đảm bảo tốc độ suy luận nhanh vừa duy trì hiệu suất phân lớp tốt, phù hợp cho hệ thống giám sát mạng xã hội theo thời gian thực.

- Triển khai thử nghiệm đối sánh hiệu năng giữa hệ thống hiện tại với các hướng tiếp cận dựa trên Prompt Engineering (Few-shot/Zero-shot learning) từ những mô hình sinh tiên tiến, nhằm đánh giá khả năng thay thế hoặc kết hợp trong tương lai.
